\section{Conclusion and Identification of the Research Gap}



The literature presents a clear picture:
\begin{enumerate}
    \item \textbf{Telepresence robots} are effective for remote diagnosis but are limited by manual control and an impersonal user experience.
    \item \textbf{Socially assistive robots} have proven their value in patient engagement and comfort but lack clinical telepresence and advanced mobility.
    \item \textbf{Autonomous mobile robots} have mastered navigation in hospitals but are used for logistics, not patient interaction.
    \item \textbf{Conversational AI} offers intelligent dialogue but lacks the physical embodiment needed for meaningful healthcare encounters.
    \item \textbf{Consumer and industrial robotics} have demonstrated that autonomous navigation technology is mature, reliable, and cost-effective for real-world deployment, as evidenced by millions of robotic vacuum cleaners in homes and over one million AMRs in Amazon warehouses.
    \item \textbf{Current hospital AMRs} are predominantly focused on material transport and disinfection, with limited integration of social interaction, advanced telepresence, and conversational AI capabilities.
\end{enumerate}

The critical research gap lies at the intersection of these domains. There is a demonstrable need for a single, integrated platform that combines the \textbf{clinical access} of a telepresence robot, the \textbf{patient-centric engagement} of a SAR, the \textbf{operational independence} of an autonomous navigator, and the \textbf{interactive intelligence} of modern conversational AI. The proven success of autonomous navigation in consumer (robotic vacuums) and industrial (Amazon warehouses) applications demonstrates that the core technology is mature and ready for advanced healthcare applications.

This project directly addresses this gap by designing a holistic system that synthesizes these strengths, creating a smart healthcare robot capable of not only connecting doctors and patients but also navigating its environment independently and engaging with patients in a helpful and empathetic manner. By learning from the successes of consumer robotics, warehouse automation, and existing hospital AMR deployments, this project aims to advance the state-of-the-art in healthcare robotics toward truly integrated, intelligent, and patient-centered care.

\input{bib/references}
